\chapter{复现命令与设备说明}
\label{app:repro}

\section{安全警告}

\texttt{contract*}、\texttt{test-recovery*}、\texttt{bench-sweep}、部分 \texttt{test-xv6.py} 目标会 \texttt{make clean} 并重建 \texttt{fs.img}；boot 自测与 \texttt{pmformat --force} 会破坏介质。请仅在可丢弃 checkout 运行并先备份。

\section{推荐验收命令}

\begin{lstlisting}
make test-model
CPUS=3 make test-vfs
make test-recovery CPUS=3 RECOVERY_ARGS="--log-dir .ci-logs/recovery"
make test-recovery-zicbom CPUS=3 RECOVERY_ARGS="--log-dir .ci-logs/zicbom"
# 补丁 QEMU 后：
# python3 tools/check_nvdimm_device.py --qemu /path/to/qemu-system-riscv64
# make test-recovery-nvdimm QEMU=/path/to/qemu-system-riscv64
make contract-full          # 破坏性旧自测
python3 tools/bench_perf.py --runs 10 --warmup 1 --samples 32 \
  --output /tmp/pmfs-matched
\end{lstlisting}

文档索引见仓库 \texttt{docs/README.md}（含语义、契约、QEMU/NVDIMM 入口）。

\section{关于「NVDIMM」与「掉电」}

\begin{itemize}[leftmargin=2em]
  \item 未打补丁的 MAP / 整机共享 RAM \textbf{不是} NVDIMM 设备。
  \item 实验补丁提供独立 \texttt{nvdimm} 冷插拔与私有 FDT，可经 QMP 核验；非上游通用、无 NFIT/热插拔。
  \item 默认 \texttt{pmem\_flush} 非硬件写回；可选 CBO 在 TCG 下仍无真实持久域。
  \item 进程 \texttt{SIGKILL} 验证的是该故障模型下的读回，\textbf{不是}物理掉电。
\end{itemize}

\section{论文 PDF 编译}

\begin{lstlisting}
cd thesis/latex
make          # xelatex + bibtex + xelatex x2
\end{lstlisting}

封面姓名/学院等宏在 \texttt{main.tex}；请替换为学校正式模板后再印刷提交。
